%%%%%%%%%%%%%%%%%%%%%%%%%%%%%%%%%%%%%%%%%%%%%%%%%%%%%%%%%%%%%%%%%%%%
%%%%%                         HEADER                           %%%%%
%%%%%%%%%%%%%%%%%%%%%%%%%%%%%%%%%%%%%%%%%%%%%%%%%%%%%%%%%%%%%%%%%%%%
%%%%%                Author : Coraline Marie                   %%%%%
%%%%%%%%%%%%%%%%%%%%%%%%%%%%%%%%%%%%%%%%%%%%%%%%%%%%%%%%%%%%%%%%%%%%


\documentclass[a4paper]{article}


%%%%% Packages %%%%%

	%%%%% Langage %%%%% 
\usepackage[frenchb]{babel}
\usepackage[utf8]{inputenc}
\usepackage[T1]{fontenc}
	%%%%% Graphique %%%%%
\usepackage{graphicx}
\usepackage{wrapfig}
	%%%%% Mise en page %%%%%
\usepackage{hyperref}
\usepackage{fancyhdr}
\usepackage{colortbl}


%%%%% Macros %%%%
%%%%%%%%%%%%%%%%%%%%%%%%%%%%%%%%%%%%%%%%%%%%%%%%%%%%%%%%%%%%%%%%%%%%
%%%%%                         MACROS                           %%%%%
%%%%%%%%%%%%%%%%%%%%%%%%%%%%%%%%%%%%%%%%%%%%%%%%%%%%%%%%%%%%%%%%%%%%
%%%%%                Author : Coraline Marie                   %%%%%
%%%%%%%%%%%%%%%%%%%%%%%%%%%%%%%%%%%%%%%%%%%%%%%%%%%%%%%%%%%%%%%%%%%%


%%%%% Ligne de séparation %%%%%
\newcommand{\HRule}{\rule{\linewidth}{0.5mm}}


%%%%% Recoloration des liens %%%%%
\hypersetup{colorlinks = true,
            breaklinks = true,
            linkcolor = {blue},
            urlcolor = {blue},}


%%%%% Création de nouvelles couleurs %%%%%
\definecolor{lightgray}{rgb}{0.75,0.75,0.75}


%%%%% Saut de ligne %%%%%
\newcommand{\NewLine}{\vspace{0.5cm}}


%%%%% Inclusion d'en-tête %%%%%
\pagestyle{fancy}
\setlength{\headheight}{23pt}	% Correction du warning \headheight is too small
\lhead{Coraline \textsc{Marie}}
\rhead{Classification de chansons}


%%%%%%%%%%%%%%%%%%%%%%%%%%%%%%%%%%%%%%%%%%%%%%%%%%%%%%%%%%%%%%%%%%%%
%%%%% DOCUMENT %%%%%
%%%%%%%%%%%%%%%%%%%%%%%%%%%%%%%%%%%%%%%%%%%%%%%%%%%%%%%%%%%%%%%%%%%%

\begin{document}

	%%%%% Page de garde %%%%%
	%%%%%%%%%%%%%%%%%%%%%%%%%%%%%%%%%%%%%%%%%%%%%%%%%%%%%%%%%%%%%%%%%%%%
%%%%%                    PAGE DE GARDE                         %%%%%
%%%%%%%%%%%%%%%%%%%%%%%%%%%%%%%%%%%%%%%%%%%%%%%%%%%%%%%%%%%%%%%%%%%%
%%%%%                Author : Coraline Marie                   %%%%%
%%%%%%%%%%%%%%%%%%%%%%%%%%%%%%%%%%%%%%%%%%%%%%%%%%%%%%%%%%%%%%%%%%%%

\begin{titlepage}
	\begin{center}

		%%% Logo et sous-titre
		\includegraphics[width=0.35\textwidth]{Figures/ATAL.png}~

		\LARGE{Master 2 \textsc{ATAL}}\\[1.5cm]

		\Large{Fouille de textes et recherche d'informations}\\[0.5cm]

		%%% Titre
		\HRule \\[0.4cm]
		{ \huge \bfseries Classification de chansons à l'aide de motifs émergents\\[0.4cm] }
		\HRule \\[1.5cm]

		%%% Auteurs et professeur
		\normalsize		
		\emph{\'Etudiante :}\\
		Coraline \textsc{Marie}

		\vspace{0.5cm}

		\emph{Encadrant :} \\
		Solen \textsc{Quiniou}

		\vspace{1cm}

		%%% Date de rendu
		{\large 12 janvier 2015}

		\vfill

		
		\includegraphics[width=0.35\textwidth]{Figures/logoUN.png}~\\[2cm]

		

	\end{center}
\end{titlepage}


	%%%%% Introduction %%%%%
	\section*{Introduction}
	\addcontentsline{toc}{section}{\protect\numberline{}Introduction}
	
	La fouille de données est une application au coeur de l'actualité. En effet, de plus en plus de données textuelles sont produites chaque jour, et sont le plus souvent non structurées et non ordonnées. C'est là tout l'intéret des applications de RI, qui permettent de sélectionner et d'extraire des informations pertinentes, dans des flux de données en éternelle extension.\\
	
	Dans le cadre du cours \textit{Fouille de textes et recherche d’informations}, dispensé au Master 2 ATAL de l'Université de Nantes, nous avions pour objectif d'utiliser des techniques de RI sur des corpus constitués de chansons. Pour cela, nous devions d'abord construire et prétraiter des corpus bien précis, puis coder un programme dont le but est d'essayer de déterminer à quel artiste appartient une chanson.\\
	
	Dans ce rapport, nous détaillerons les différents choix et démarches adoptées pour ce projet. Nous étudierons donc d'abord les artistes et la construction des corpus. Puis nous verrons quels ont été les prétraitements et les méthodes utilisées pour permettre l'extraction de motifs émergents, avant de finir sur les résultats produits par le programme.

	%%%%% Parties à ajouter %%%%%
	\section{Construction des corpus}
	
	\subsection{Choix des artistes}
	
	De nos jours, la musique est présente partout, dans toutes les langues et dans de nombreux styles musicaux différents. La première démarche de ce projet fût donc de choisir trois artistes, dont les oeuvres soient à la fois nombreuses, et dans la même langue. \\
	
	Le premier artiste que j'ai choisi est Linkin Park. Ce groupe a déjà produit une quinzaine d'albums, tous en anglais, dont la plupart dans un style proche du hard rock, et du néo métal. J'ai donc retenu quatre albums : The Hunting Party (2014), Hybrid Theory (2000), Meteora (2003) et Living Things (2012).\\
	
	Le second artiste que j'ai sélectionné ressemble relativement à Linkin Park de part son style musical (hard rock et néo métal), son style de paroles et sa langue. Il s'agit du groupe Breaking Benjamin, auteur de sept albums, dont quatre utilisés pour ce projet : Dear Agony (2009), We Are Not Alone (2004), Phobia (2006) et Saturate (2002).\\
	
	Le dernier artiste que j'ai choisi diffère complètement des deux précédents : Daft Punk. Bien que ce groupe français produise des oeuvres en Anglais, leur style musical est plus électronique et groovy. Pour cet artiste, j'ai sélectionné les quatre albums suivants : Random Access Memories (2013), Homework (1997), Discovery (2001) et Human after all (2005).
	
	\subsection{Création de corpus}
	
	Pour chaque artiste, j'ai créé un corpus d'apprentissage, et un corpus de test. Les corpus d'apprentissage serviront à l'entraînement des données, alors que les corpus de test seront utilisés pour essayer de trouver à quels artistes ils appartiennent.
	
	\subsubsection{Corpus d'entraînement}
	
	Pour ce projet, trois corpus d'entraînement ont été créés : un pour chaque artiste. Ces corpus sont contitués de deux albums et demi par artiste, ce qui signifie que deux albums entiers et la moitié d'un troisième ont été utilisés pour construire ces corpus. Les albums n'ont pas été sélectionné aléatoirement, car à chaque fois, il y a un album récent et plus ancien dans le corpus d'apprentissage. Ce choix a été fait intentionnellement pour que l'impact de l'évolution de langage d'un artiste dans le temps, ait le moins d'influence possible sur les motifs.\\
	
	Ainsi, le corpus d'apprentissage de Linkin Park contient les albums "The Hunting Party" (2014), "Hybrid Theory" (2000), et la moitié des chansons de "Meteora" (2003). Ceci représente au total 28 chansons, soit 1535 vers pour un peu moins de 10 000 mots.\\
	
	Le corpus d'entrainement de Breaking Benjamin est composé des albums "Dear Agony" (2009), "We Are Not Alone" (2004) et de la moitié de "Phobia" (2006). Ces chansons sont au nombre de 28, ce qui fait 1204 vers, pour environ 6000 mots.\\
	
	Pour finir, le corpus d'apprentissage de Daft Punk a été formé à partir des albums "Random Access Memories" (2013), "Homework" (1997) et la moitié de "Discovery" (2001). Ce corpus contient 21 chansons, c'est à dire 877 vers et environ 4800 mots.
	
	\subsubsection{Corpus de test}
	
	Les corpus de tests créés sont légèrement différents des corpus d'apprentissage. En effet, ceux-ci doivent être analysés par le programme qui doit déterminer ensuite à quel artiste appartient chaque chanson. Ainsi, il existe un corpus de test par artiste, et chacun contient un fichier par chanson. \\
	
	Le corpus de test de Linkin Park contient la moitié des chansons de l'album "Meteora" (2003), et toutes les chansons de "Living Things" (2012). Ce qui représente au total 14 chansons de 567 vers et environ 3700 mots.\\
	
	En ce qui concerne le corpus de test de Breaking Benjamin, celui-ci contient la moitié de l'album "Phobia" (2006) et toutes les chansons de "Saturate" (2002). Ce qui donne 18 chansons de 659 vers et d'environ 3350 mots.\\
	
	Le corpus de test de Daft Punk est quant à lui composé de la moitié des chansons de l'album "Discovery" (2001) et de toutes celles de "Human after all" (2005). On dénombre donc 9 chansons avec 343 vers et environ 2100 mots.
	
	
	\section{Extraction de motifs}
	
	\subsection{Prétraitement des corpus}
	
	Après avoir fini de construire les corpus, nous devions extraire un certain nombre de motifs émergents. Cependant, les corpus sont extrèmement bruités, il a donc fallu au préalable les prétraiter. Pour cela, j'ai commencé par transformer toutes les majuscules en minuscules, et tous les éléments de ponctuations comme les '!' et les '?' en '.'. En analysant les corpus, j'ai également remarqué que certaines phrases sont composées de plusieurs vers, mais qu'aucune ponctuation n'y a été placée. J'ai donc rajouté des '.' aux endroits qu'il me paraissait être le mieux, pour donner un sens aux groupes de vers. J'ai également supprimé les smileys, et les précisions de chants (chorus, single, x2, \dots).\\
	
	Dans un second temps, j'ai également voulu supprimer les stopwords, cependant je ne l'ai finalement pas fait pour deux raisons. Tout d'abord car cela diminue énormément la taille de mes corpus (plus de 80\% des mots sont supprimés dans le corpus d'apprentissage de Linkin Park), et ensuite car je pense que les stopwords peuvent être représentatif du style d'un artiste. En effet, certains artistes emploient beaucoup plus de stopwords que d'autres, et cela peut être un moyen de différencier le style de parole lors de la reconnaissance des chansons.
	
	\subsection{Extraction de motifs automatique}
	
	Une fois que tous les corpus ont été prétraités, il faut en extraire des motifs. Pour cela, je me suis servie de Greyc sur le site SDMC\footnote{https://sdmc.greyc.fr}. Sur cet outil, plusieurs paramètres peuvent être modifiés pour l'extraction de motif automatique. Pour ma part, j'ai utilisé 8 combinaisons différentes de paramètres :
	\begin{table}[htb]
          \centering
          \begin{tabular}{cccc}
          \noalign{\smallskip} \hline \hline \noalign{\smallskip}
          Config & gap & taille & support\\
          \hline \noalign{\smallskip}
          \no 1 & [0;5] & 10 & 10 \\ 
          \no 2 & [0;0] & 10 & 10 \\ 
          \no 3 & [0;5] & 10 & 20 \\ 
          \no 4 & [0;0] & 10 & 20 \\ 
          \no 5 & [0;5] & 10 & 5 \\ 
          \no 6 & [0;0] & 10 & 5 \\ 
          \no 7 & [0;5] & 10 & 3 \\ 
          \no 8 & [0;0] & 10 & 3 \\
          \noalign{\smallskip} \hline \noalign{\smallskip}
          \end{tabular}
          \caption{Paramètres utilisés sur Greyc}
     \end{table}
	
	J'ai également voulu faire varier le paramètre concernant la taille maximale du motif. Cependant la majeur partie de mes motifs sont au maximum de taille 3, donc faire varier ce paramètre n'influe quasiment pas sur le nombre de motifs produit.\\
	
	\begin{table}[htb]
          \centering
          \begin{tabular}{cccc}
          \noalign{\smallskip} \hline \hline \noalign{\smallskip}
          Config & Linkin Park & Breaking Benjamin & Daft Punk\\
          \hline \noalign{\smallskip}
          \no 1 & 633 & 284 & 239 \\ 
          \no 2 & 297 & 169 & 164 \\ 
          \no 3 & 197 & 102 & 97 \\ 
          \no 4 & 117 & 74 & 74 \\ 
          \no 5 & 1813 & 3777 & 452 \\ 
          \no 6 & 654 & 385 & 271 \\ 
          \no 7 & 10155 & 14257 & 1123 \\ 
          \no 8 & 1002 & 713 & 373 \\
          \noalign{\smallskip} \hline \noalign{\smallskip}
          \end{tabular}
          \caption{Nombre de motifs obtenu en fonction des paramètres choisis}
     \end{table}
	
	En ce qui concerne les corpus de tests, j'ai utilisé les mêmes paramètres à l'exception du dernier, car dans une chanson d'en moyenne 30 vers, il est difficile d'obtenir des motifs qui apparaissent au moins 20 fois.
	
	\subsection{Extraction de motifs émergents}
	
	L'extraction de motifs automatique par l'intermédiaire de Greyc, m'a permis d'obtenir de nouveaux fichiers d'entraînement. Cependant, certains motifs sont communs à plusieurs de ces fichiers, il faut donc les supprimer. C'est pourquoi nous avons créé trois "classifieurs" différents. Le premier va permettre de comparer Linkin Park à Breaking Benjamin, le second permet de comparer Linkin Park à Daft Punk, et le dernier compare Linkin Park au deux autres artistes.\\
	
	Ainsi, pour éviter au mieux au programme toute confusion entre les artistes, nous avons extrait les motifs émergents de chaque artiste en fonction du classifieur où il se trouve. Pour ce faire, j'ai codé un script qui compare tous les motifs des artistes et qui ne conserve que les motifs émergents.
	
	\section{Résultats}
	
	\subsection{\'Evaluation}
	
	L'évaluation se fait par l'intermédiaire d'un script qui évalue la précision, le rappel et la f-mesure pour chaque classifieur. Bien que les résultats ne soient pas très satisfaisant, je vais les détailler ci-après. Chaque tableau présente les valeurs calculées pour tous les artistes appartenant au même classifieur et pour toutes les configurations de paramètres testées avec Greyc. 
	
	\begin{table}[!htb]
          \centering
          \begin{tabular}{ccccccc}
          \noalign{\smallskip} \hline \hline \noalign{\smallskip}
          Config & \multicolumn{3}{c}{Linkin Park} & \multicolumn{3}{c}{Breaking Benjamin} \\
          ~ & P & R & F & P & R & F \\
          \hline \noalign{\smallskip}
          \no 1 & 0.7 & 0.5 & 0.58 & 0.87 & 0.39 & 0.54\\ 
          \no 2 & 0.67 & 0.43 & 0.52 & 0.78 & 0.39 & 0.52\\ 
          \no 3 & 0.54 & 0.5 & 0.52 & 0.8 & 0.22 & 0.35\\ 
          \no 4 & 0.6 & 0.43 & 0.5 & 0.86 & 0.33 & 0.48\\ 
          \no 5 & 0.63 & 0.36 & 0.46 & 0.25 & 0.22 & 0.24\\ 
          \no 6 & 0.6 & 0.43 & 0.5 & 0.8 & 0.22 & 0.35\\ 
          \no 7 & 0 & 0 & 0 & 0.56 & 1 & 0.72\\ 
          \no 8 & 0.54 & 0.43 & 0.48 & 1 & 0.11 & 0.2\\
          \noalign{\smallskip} \hline \noalign{\smallskip}
          \end{tabular}
          \caption{Résultats obtenus pour le premier classifieur en fonction de la configuration des paramètres}
	
          \centering
          \begin{tabular}{ccccccc}
          \noalign{\smallskip} \hline \hline \noalign{\smallskip}
          Config & \multicolumn{3}{c}{Linkin Park} & \multicolumn{3}{c}{Daft Punk} \\
          ~ & P & R & F & P & R & F \\
          \hline \noalign{\smallskip}
          \no 1 & 0.85 & 0.79 & 0.81 & 0 & 0 & 0\\ 
          \no 2 & 0.85 & 0.79 & 0.81 & 0 & 0 & 0\\ 
          \no 3 & 0.8 & 0.86 & 0.83 & 0 & 0 & 0\\ 
          \no 4 & 0.8 & 0.86 & 0.83 & 0 & 0 & 0\\ 
          \no 5 & 0.91 & 0.72 & 0.8 & 0 & 0 & 0\\ 
          \no 6 & 0.92 & 0.79 & 0.85 & 0 & 0 & 0\\ 
          \no 7 & 0 & 0 & 0 & 0.22 & 0.44 & 0.30\\ 
          \no 8 & 0.92 & 0.79 & 0.85 & 0 & 0 & 0\\
          \noalign{\smallskip} \hline \noalign{\smallskip}
          \end{tabular}
          \caption{Résultats obtenus pour le second classifieur en fonction de la configuration des paramètres}
	
          \centering
          \begin{tabular}{cccccccccc}
          \noalign{\smallskip} \hline \hline \noalign{\smallskip}
          Config & \multicolumn{3}{c}{Linkin Park} & \multicolumn{3}{c}{Breaking Benjamin} & \multicolumn{3}{c}{Daft Punk}\\
          ~ & P & R & F & P & R & F & P & R & F \\
          \hline \noalign{\smallskip}
          \no 1 & 0.63 & 0.36 & 0.45 & 0.67 & 0.33 & 0.44 & 0 & 0 & 0\\ 
          \no 2 & 0.63 & 0.36 & 0.45 & 0.56 & 0.28 & 0.37 & 0 & 0 & 0\\ 
          \no 3 & 0.4 & 0.43 & 0.41 & 0.6 & 0.17 & 0.26 & 0 & 0 & 0\\ 
          \no 4 & 0.46 & 0.36 & 0.4 & 0.67 & 0.22 & 0.33 & 0 & 0 & 0\\ 
          \no 5 & 0.67 & 0.29 & 0.4 & 0.57 & 0.22 & 0.32 & 0 & 0 & 0\\ 
          \no 6 & 0.56 & 0.36 & 0.43 & 0.72 & 0.28 & 0.4 & 0 & 0 & 0\\ 
          \no 7 & 1 & 0.07 & 0.13 & 0.52 & 0.61 & 0.56 & 0.33 & 0.11 & 0.17\\ 
          \no 8 & 0.56 & 0.36 & 0.44 & 0.75 & 0.17 & 0.27 & 0 & 0 & 0\\
          \noalign{\smallskip} \hline \noalign{\smallskip}
          \end{tabular}
          \caption{Résultats obtenus pour le troisième classifieur en fonction de la configuration des paramètres}
	\end{table}
	
	\subsection{Discussion}
	
	En regardant attentivement les résultats fournis par les tables 3, 4 et 5, plusieurs remarques peuvent être faites. Tout d'abord, rappelons que le premier classifieur comparait les chansons des deux artistes qui se ressemblent le plus. Ainsi, on remarque que pour la majorité des configurations, le programme attribue environ 50\% des chansons de Linkin Park au bon artiste, et seulement 30\% celles de Breaking Benjamin. Ces erreurs sont sûrement dûes au fait que les artistes sont assez proches au niveau de leurs paroles, car ils ont beaucoup de motifs en commun. De plus, la structure des paroles et le prétraitement réalisé ne sont peut-être pas appropriés pour ce travail.\\
	
	En ce qui concerne le second classifieur, les résultats sont des plus décevants. En effet, étant donné que Daft Punk et Linkin Park sont des artistes très différents, on aurait pu imaginer que le programme puissent aisément différencier les deux artistes, mais ce n'est pas le cas. Si on regarde les résultats pour Linkin Park, le programme reconnaît plutôt bien ses chansons, cependant il est incapable de reconnaître celles de Daft Punk. En regardant plus attentivement les causes de ces échecs, on remarque que les corpus de Daft Punk sont beaucoup moins fournis que ceux de Linkin Park. Ainsi, il y a moins de données d'apprentissage, et donc moins de motifs émergents pertinents. De plus, Daft Punk est un artiste qui utilise énormément la répétition dans ses chansons, c'est pourquoi il y a encore moins de motifs, car il y a moins de mots différents. \\
	
	Le dernier classifieur regroupe les trois artistes. Ainsi, comme dans le premier classifieur, on remarque qu'environ 40\% des chansons de linkin Park sont reconnues, et moins de 30\% de celles de Breaking Benjamin. Cette légère baisse par rapport au premier classifieur peut s'expliquer car la présence moindre de motifs émergents dans ce dernier classifieur. En ce qui concerne Daft Punk, les résultats sont toujours décevants, car le programme ne reconnaît toujours pas ces chansons.\\
	
	En outre, nous pouvons également remarquer certains résultats surprenants, comme ceux de la 7ème configuration, car contrairement aux autres, les chansons de Linkin Park ne sont pas reconnues alors que certaines de Daft Punk le sont.
	
	\section{Conclusion}
	
	Au terme de ce projet, je suis assez satisfaite de ce travail. En effet, contrairement à ce que l'on a l'habitude de faire en cours, ce projet nécessitait des données plus agréables à travailler. De plus, il nous a permis de mettre en pratique plusieurs notions étudiées en cours, et pas seulement en "Fouille de textes et recherche d'informations". Il a donc fallu reprendre des concepts étudiés en "corpus", en "Alignement de chaînes et de textes" et en "Machine Learning". Ainsi, malgré les résultats décevants que j'ai obtenu, je suis tout de même satisfaite car ce projet m'a permis de tester et d'améliorer mes compétences.

	%\newpage
	
	%%%%% Bibliographie %%%%% (Attention, citer une source de bibliographie sinon bug)
	%\bibliographystyle{plain}
	%\bibliography{biblio}

\end{document}
